\documentclass[twocolumn]{aastex701}

\usepackage{amsmath}
\usepackage{amssymb}

\newcommand{\harmonica}{\texttt{harmonica}}
\newcommand{\jax}{\textsc{jax}}
\newcommand{\dd}{\mathrm{d}}

\shorttitle{Harmonica on GPUs}
\shortauthors{Fairnington}

\begin{document}

\title{Harmonica v2: Fixed-Work Transmission-String Light Curves for GPUs}

\author{Tyler Fairnington}
\affiliation{Department of Astronomy and Astrophysics, University of
Chicago, Chicago, IL 60637, USA}
\email{tfairnington@uchicago.edu}

\begin{abstract}
\harmonica\ computes transit light curves of planets whose sky-projected
silhouette is a Fourier series in angle, a ``transmission string''. Its
C++ core finds the planet--star limb intersections from the eigenvalues
of a companion matrix and integrates the occulted intensity with Green's
theorem, branching on the geometry of each point. That structure is fast
on one CPU core but maps poorly onto GPUs. We describe a reformulation in
which every light-curve point performs the same, statically shaped work,
so that the model compiles with XLA, runs on GPUs, batches over many
light curves, and composes with automatic differentiation, including to
higher order. Intersections are bracketed on angular grids, with a closed
form for circles and a proven two-extremum bracket for first-order
strings, and their derivatives follow from the implicit-function theorem.
The occulted region is assembled from the direction in which the limbs
cross, which we show handles tangencies and stellar arcs that wrap around
the star. Validated against independent brute-force overlap areas and
analytic cases, the model agrees with the C++ to ${\sim}10^{-10}$ at the
default quadrature order and to ${\sim}4\times10^{-15}$ at matched order,
where the C++ is itself correct. Validation also uncovered several errors
in the original implementation, including fluxes above unity for planets
that cover the stellar centre and a flattened orbit near apastron. A
dispatch layer evaluates the C++ on CPUs, with correct higher
derivatives, and the new model elsewhere.
\end{abstract}

\keywords{Exoplanet atmospheres (487) --- Transits (1711) ---
Astronomy software (1855) --- GPU computing (1969)}

\section{Introduction} \label{sec:intro}

Transmission strings describe the radius of a transiting planet as a
function of sky-plane angle $\theta$ about its centre,
\begin{equation}
r_\mathrm{p}(\theta) = a_0 + \sum_{n=1}^{N_c} \left[a_n \cos n\theta +
b_n \sin n\theta\right],
\label{eq:string}
\end{equation}
so that the light curve constrains how the atmospheric opacity varies
around the terminator \citep{grant2023}. The order $N_c$ sets the angular
resolution: $N_c=0$ is a circle, while $N_c=1$ adds a first-order
asymmetry, for example between the morning and evening limbs. Recent
JWST analyses use strings of order six \citep{alderson2026}.

Inference with these models increasingly involves many light curves at
once, e.g.\ hundreds of spectroscopic channels, cross-validation folds or
ensemble samplers. GPUs suit this workload, but only for programs whose
control flow and array shapes do not depend on the data. The original
\harmonica\ evaluates each point with a data-dependent sequence of
branches, an eigenvalue problem whose number of real roots varies, and a
quadrature order that switches with position. Wrapping it for a GPU would
therefore not help: the algorithm itself must change. This Letter
describes those changes, released on the \texttt{jax-gpu} branch of
\harmonica.\footnote{\url{https://github.com/TylerFair/harmonica/tree/jax-gpu}}

\section{The model} \label{sec:model}

We follow the formulation of \citet{grant2023}. For limb-darkening laws
that are sums of powers of $\mu$, $I(\mu) = \sum_n p_n \mu^n$, the
normalised flux is
\begin{equation}
F = 1 - I_0 \sum_n p_n\, s_n,
\end{equation}
where $s_n$ is the occulted integral of $\mu^n$. By Green's theorem, each
$s_n$ is a line integral around the boundary of the occulted region,
which consists of arcs of the planet's limb and of the stellar limb. In
planet-centred polar coordinates, with the planet at separation $d$ and
angle $\nu$ from the stellar centre, a planet-limb arc contributes
\begin{equation}
s_n^{(p)} = \int_{\theta_j}^{\theta_{j+1}} \zeta_n(z_\mathrm{p})\,
\eta(\theta)\, \dd\theta,
\label{eq:planet_arc}
\end{equation}
with
\begin{align}
\eta &= r_\mathrm{p}^2 - d\, r_\mathrm{p} \cos(\theta-\nu)
       - d\, r_\mathrm{p}' \sin(\theta-\nu), \\
z_\mathrm{p}^2 &= 1 - d^2 - r_\mathrm{p}^2
       + 2 d\, r_\mathrm{p} \cos(\theta-\nu), \\
\zeta_n(z) &= \frac{1 - z^{n+2}}{(n+2)(1-z^2)}.
\label{eq:zeta}
\end{align}
Here $z_\mathrm{p}$ is $\mu$ on the planet's limb. A stellar-limb arc,
on which $\mu=0$, contributes $\Delta\phi/(n+2)$, where $\Delta\phi$ is
the arc's angular extent about the stellar centre.

The quadratic law has $p = (1-u_1-u_2,\, u_1+2u_2,\, -u_2)$ for
$\mu^{0,1,2}$ and $I_0^{-1} = \pi(1 - u_1/3 - u_2/6)$. The non-linear law
has $p = (1-\sum_i u_i,\, u_1, u_2, u_3, u_4)$ for $\mu^{0,1/2,1,3/2,2}$
and $I_0^{-1} = \pi(1 - u_1/5 - u_2/3 - 3u_3/7 - u_4/2)$. We add the
power-2 law $I = 1 - c(1-\mu^\alpha)$ \citep{hestroffer1997, maxted2018},
with $p=(1-c,\, c)$ for $\mu^{0,\alpha}$ and
$I_0^{-1} = \pi[1 - c + 2c/(\alpha+2)]$. With $\alpha=2$ it reduces
exactly to the quadratic law with $u_1=2c$, $u_2=-c$.

\section{A fixed-work formulation} \label{sec:method}

\subsection{Limb intersections} \label{sec:roots}

The limbs intersect where
\begin{equation}
f(\theta) \equiv r_\mathrm{p}^2 + d^2 - 2 d\, r_\mathrm{p}
\cos(\theta-\nu) - 1 = 0,
\label{eq:f}
\end{equation}
with $f<0$ inside the star. Since $r_\mathrm{p}$ is a trigonometric
polynomial of degree $N_c$, $f$ is one of degree $D = \max(2N_c, N_c+1)$,
and has at most $2D$ real roots on the circle. The C++ finds them as
eigenvalues of a $2D\times2D$ companion matrix, which requires $a_{N_c}$
and $b_{N_c}$ not both to vanish. Instead we use three bracketing schemes,
each followed by bisection.

\textit{Circles.} When every $a_n, b_n$ ($n\geq1$) is zero, the roots are
$\nu\pm\beta$, with
\begin{equation}
\beta = \operatorname{atan2}\!\left(\sqrt{\Delta},\; d^2 + a_0^2 - 1\right),
\end{equation}
where $\Delta = [1-(d-a_0)^2][(d+a_0)^2-1]$. The factored discriminant
stays accurate at contact, where the usual $\arccos$ form loses half its
digits.

\textit{First-order strings} ($N_c=1$). Here $f$ contains only the first
and second harmonics,
\begin{equation}
f = c_0 + A\cos(\theta-\phi) + B\cos(2\theta-\psi),
\end{equation}
with $A\cos\phi = 2a_0(a_1 - d\cos\nu)$, $A\sin\phi = 2a_0(b_1 - d\sin\nu)$,
$B\cos\psi = (a_1^2-b_1^2)/2 - d(a_1\cos\nu - b_1\sin\nu)$ and
$B\sin\psi = a_1 b_1 - d(a_1\sin\nu + b_1\cos\nu)$.

\textit{Proposition.} If $A > \sqrt{20}\,B$, then $f$ has exactly one
maximum in $(\phi-\pi/2, \phi+\pi/2)$ and one minimum in
$(\phi+\pi/2, \phi+3\pi/2)$, and no other stationary points.

\textit{Proof.} Let $x = \theta-\phi$. Then
$f' = -A\sin x - 2B\sin(2\theta-\psi)$ vanishes only where
$|\sin x| \leq 2B/A$, and there
$|\cos x| \geq (1-4B^2/A^2)^{1/2}$. Hence on that set
$|f''| \geq A(1-4B^2/A^2)^{1/2} - 4B$, which is positive exactly when
$A^2 > 20B^2$. So $f'$ is strictly monotone on each of the two arcs on
which it can vanish. Since $2B < A$, $f'(\phi-\pi/2) > 0 > f'(\phi+\pi/2)$,
and similarly with the signs exchanged at $\phi+\pi/2$ and $\phi+3\pi/2$,
so each arc contains exactly one zero. $\square$

When the condition holds, which it does for any physically small
asymmetry, both extrema are bisected on these certified brackets. Roots
exist if and only if $f$ is positive at the maximum and negative at the
minimum, and one lies on each arc between them, on which $f$ is monotone.
This finds a pair of intersections arbitrarily close to contact, which a
fixed angular grid can miss entirely.

\textit{General strings.} Otherwise, including uncertified $N_c=1$
shapes, we evaluate $f$ and $f'$ on a uniform grid of
$G=\max(128, 64N_c)$ intervals of width $h=2\pi/G$, bisect every sign
change of $f'$, and insert the resulting extrema into the grid. Sign
changes of $f$ between consecutive nodes then bracket the roots. The
nodes form a closed loop. A node at which $f$ is exactly zero takes the
sign of the preceding node, so that a crossing on a node is counted once
and a tangency, which does not change sign, is not counted at all. A
closed curve then crosses an even number of times, at most $2D$; any
other count means the bracketing failed, and the flux is returned as NaN
rather than as a plausible but wrong value.

This scheme is not certified. It finds every transverse root if each grid
interval contains at most one zero of $f'$, e.g.\ if distinct zeros of
$f'$ are separated by more than $h$. If $|f'''|\leq L$ and
$|f''|\geq m>0$ at every critical point, Taylor's theorem gives a
separation of at least $2m/L$, so $h < 2m/L$ suffices. A computable
per-cell alternative, with midpoint $c$ and half-width $q=h/2$, is that
either $|f''(c)| > Lq$, so that $f'$ is monotone on the cell, or
$|f'(c)| > |f''(c)|\,q + Lq^2/2$, so that it has no zero there. Writing
$A_n = (a_n^2+b_n^2)^{1/2}$, $R_0 = |a_0| + \sum_n A_n$ and
$R_j = \sum_n n^j A_n$, a conservative bound is
$L = 6R_1R_2 + 2R_0R_3 + 2|d|(R_3+3R_2+3R_1+R_0)$. A cell failing both
tests is uncertified, not necessarily wrong.

\subsection{Root derivatives} \label{sec:jvp}

Differentiating through bisection gives derivatives of zero. Instead, the
root finder is given a custom derivative rule from the implicit-function
theorem, $\dot\theta_j = -\dot f/f_\theta$, with
\begin{align}
f_\theta &= 2 r_\mathrm{p} r_\mathrm{p}'
  - 2d\left(r_\mathrm{p}' C - r_\mathrm{p} S\right),\\
\dot f &= 2\left(r_\mathrm{p} - dC\right)\dot r_\mathrm{p}
  + 2\left(d - r_\mathrm{p} C\right)\dot d
  - 2 d\, r_\mathrm{p} S\,\dot\nu,
\end{align}
with $C = \cos(\theta-\nu)$ and $S = \sin(\theta-\nu)$, and
$\dot r_\mathrm{p}$ is Equation~\eqref{eq:string} evaluated with the
coefficient tangents. Padded roots, and roots with
$|f_\theta| < 10^{-14}$, receive zero tangents. The denominator is made
safe \emph{before} dividing, since masking a $0/0$ afterwards still leaks
NaN into reverse-mode derivatives. Everything downstream of the roots is
differentiated automatically.

\subsection{Assembling the occulted region} \label{sec:assembly}

The C++ classifies each pair of consecutive crossings with a lookup table
of crossing types, and integrates the stellar arc between them using the
principal difference of their stellar angles. We found that this
mis-assembles some configurations (Section~\ref{sec:issues}), and instead
build the boundary from the direction in which the limbs cross.

\textit{Planet arcs.} Order the crossings by $\theta$. Since $f<0$
inside the star, the planet-limb arc after a crossing bounds the occulted
region exactly when $f_\theta<0$ there.

\textit{Stellar arcs.} Order the crossings instead by their angle $\phi$
about the stellar centre, and take each arc anticlockwise to the next,
wrapping the last by $2\pi$. The arc after a crossing bounds the region
exactly when $f_\theta > 0$ there. To see this, parametrise the stellar
limb by $\phi$, with $\rho(\phi)$ and $\theta(\phi)$ the distance and
angle about the planet's centre, and let
$H(\phi) = r_\mathrm{p}^2(\theta(\phi)) - \rho^2(\phi)$, which is positive
inside the planet. At a crossing, $\theta_\phi = (1-d\cos\phi)/r_\mathrm{p}^2$,
$\rho_\phi = d\sin\phi/r_\mathrm{p}$, and
$1-d^2 = r_\mathrm{p}^2-2d\,r_\mathrm{p}\cos(\theta-\nu)$, from which
$H_\phi = f_\theta$. $H$ cannot change sign before the next transverse
crossing, and tangencies omitted from the root list do not change it.

This requires positive radii and a complete list of transverse
crossings, but no midpoint tests, crossing-type tolerances or angle
wrapping heuristics. Every point evaluates $K=2D$ potential arcs of each
kind, masked by these conditions, and masked arcs are given benign
geometry so that they cannot produce non-finite derivatives.

Configurations without crossings are resolved in a fixed order. Bounds
$r_\mathrm{p}\in[a_0 - \sum_n A_n,\, a_0 + \sum_n A_n]$, with a $10^{-12}$
tolerance, identify planets wholly inside or outside the star and stars
wholly inside the planet; they are exact for $N_c\leq1$ and conservative
otherwise. Such a trivial classification overrides any roots and
integrates the entire limb, since otherwise a root pair within the
tolerance of contact would be integrated over only its tiny connecting
arc. With no roots otherwise, the limb point facing the stellar centre,
$r_\mathrm{p}(\nu)$, decides between the remaining cases.

\subsection{Line integrals} \label{sec:quadrature}

The C++ evaluates $s_0$ and $s_2$ in closed form and the other terms by
Gauss--Legendre quadrature, of order 20 or 50 depending on position. We
apply a single order $n_\mathrm{GL}$ (default 50) to every planet-limb
term, which removes the position-dependent branch and is at least as
accurate as the C++ default (Table~\ref{tab:accuracy}).

Equation~\eqref{eq:zeta} has a removable singularity at $z=1$, the
stellar disc centre, and cancels catastrophically near it. We use exact
rearrangements that avoid both:
\begin{align}
\zeta_0 &= \tfrac{1}{2}, \qquad
\zeta_2 = \tfrac{1}{4}(1+z^2), \\
\zeta_1 &= \tfrac{1}{3}\left[z + (1+z)^{-1}\right], \\
\zeta_{1/2} &= \frac{2}{5}\,\frac{\sum_{k=0}^{4} s^k}{(1+s)(1+s^2)}, \\
\zeta_{3/2} &= \frac{2}{7}\,\frac{\sum_{k=0}^{6} s^k}{(1+s)(1+s^2)},
\end{align}
where $s = z^{1/2}$. For the power-2 term, with $t = 1-z^2$ and
$p=(\alpha+2)/2$, we sum twelve terms of the Taylor series of
$[1-(1-t)^p]/(2pt)$ when $t < \min(0.01, 0.1/p)$, which keeps its
truncation error below ${\sim}10^{-21}$ for any $\alpha$; evaluate
\begin{equation}
\zeta_\alpha = -\frac{\operatorname{expm1}\!\left[p\,
\operatorname{log1p}(-t)\right]}{2pt}
\end{equation}
for larger $t < 1/2$; and use Equation~\eqref{eq:zeta} directly beyond.
Values and first and second derivatives were checked against a 75-digit
reference for $\alpha$ up to 100, with errors below $3\times10^{-12}$;
no finite series supplies correct derivatives of arbitrarily high order.

\subsection{Orbits} \label{sec:orbit}

The orbit follows the C++, including its rotation of the string with the
orbit, but with three changes. Kepler's equation is solved on $[0,\pi]$ by
Newton iterations safeguarded by bisection, on
$(1-e)E + e(E-\sin E) = M$ with a series for $E-\sin E$ near zero, so that
it does not cancel as $e\to1$; its derivatives are given by the
implicit-function theorem. The true anomaly and radius then follow from
$s=\sin(E/2)$, as $\cos f = [(1-e)-2s^2]/[(1-e)+2e s^2]$,
$\sin f = \sqrt{(1-e)(1+e)}\,\sin E/[(1-e)+2es^2]$ and
$r=a[(1-e)+2es^2]$, which avoid both cancellation near $e=1$ and the
singular half-angle formula, $\tan(f/2) \propto \tan(E/2)$, near apastron.
Finally, the eccentric anomaly at transit is
$E_0 = \operatorname{atan2}(\sqrt{(1-e)(1+e)}\cos\omega,\, e + \sin\omega)$,
whose derivatives are regular at $\omega = -\pi/2$. When $e$ is a
compile-time zero, Kepler's equation is skipped; when it is traced, the
eccentric solution is always used, so that $\partial F/\partial e$ is
correct at $e=0$.

\subsection{Derivatives and the CPU path} \label{sec:dispatch}

On a single CPU core, padding every point to the most complex geometry
costs about an order of magnitude relative to the C++. The public
functions therefore choose an implementation per platform at compile
time (\texttt{jax.lax.platform\_dependent}): the C++ on CPUs, and the
fixed-work model elsewhere. To make this possible, we moved the C++ \jax\
primitives to \jax's public foreign function interface, as the APIs they
used have been removed from \jax. The C++ supplies only first
derivatives, so its Jacobian is given its own derivative rule, the
fixed-work model's Hessian-vector product. That rule is traced only for
second and higher derivatives: compiled first-order programs contain no
fixed-work evaluations and keep the C++'s speed. The flux depends on time
only through $t-t_0$, which gives the time derivatives on this path.

Automatic differentiation gives exact derivatives of the implemented
function. These approximate the model's derivatives to the accuracy of
the quadrature and series above, and only where the topology is
resolved. At a generic contact with a limb-darkened star, the occulted
flux grows locally as $\delta^{3/2}$ in the penetration depth $\delta$,
so the model's second derivatives are themselves singular there, and
finite derivatives assigned at contact or tangency (e.g.\ the
$|f_\theta|$ guard) do not make it smooth.

\section{Validation} \label{sec:validation}

All tests use double precision, which the model requires. Agreement with
the C++ alone cannot establish correctness (Section~\ref{sec:issues}),
so we lead with independent references.
\begin{itemize}
\item \textit{Overlap areas.} For a uniform star, $1-F$ is the overlap
area over $\pi$, which we integrate independently in closed form along
rays from the planet's centre. Over 6000 random positive-radius shapes,
$N_c=1$--4, $a_0 = 0.2$--1.2 and $d$ anywhere, including planets that
cover the stellar centre and multi-lobed overlaps, the largest
differences are ${\sim}10^{-7}$; these come from the reference itself,
and fall to ${\sim}10^{-10}$ when its resolution is raised from
$2\times10^5$ to $8\times10^6$ rays.
\item \textit{Analytic cases.} The shape $r_\mathrm{p} = 1 + a\cos\theta
- a\cos^2\theta$ centred on the star has crossings at $\pm\pi/2$ and a
tangency at $\theta=0$, and
$F = a/2 + 2a/\pi - 7a^2/16 - 4a^2/(3\pi)$; we agree to $10^{-15}$. We
also reproduce a tangency coexisting with crossings,
$r_\mathrm{p} = 0.125 - 0.00390625\cos 6\theta$ at $d=1.12109375$
($F = 0.999999975014641$), and nearby four-crossing cases.
\item \textit{C++ parity.} For $N_c = 0$--3, circular and eccentric
orbits, and both C++ laws, the default order agrees to
$\leq 3\times10^{-10}$, and matched quadrature to ${\sim}4\times10^{-15}$
(Table~\ref{tab:accuracy}), in configurations where the C++ is correct.
Power-2 agrees with an independent C++ implementation of the law to
$1.2\times10^{-10}$ for $N_c\leq2$.
\item \textit{Derivatives.} Reverse- and forward-mode derivatives agree
to $5\times10^{-14}$. They agree with the C++ analytic derivatives to
$\leq10^{-5}$, the largest differences being just inside contact, where
ours agree with converged finite differences to ${\sim}10^{-10}$
(Section~\ref{sec:issues}). Through the CPU path, the Hessian of a planet
inside a uniform star is $-2$ as required, and partial-transit Hessians
match the fixed-work model's to $5\times10^{-15}$.
\item \textit{Orbits.} The derivative of $\sin f$ at apastron is exact,
and the true anomaly at $e = 0.99999999$ near periastron matches a
long-double reference to $10^{-8}$.
\item \textit{Inference.} Re-running the \harmonica\ gradient-based
inference tutorial ($N_c=3$, NUTS via NumPyro; \citealt{phan2019})
reproduces its published posterior within sampling noise, with
$\hat{R}\leq1.008$.
\end{itemize}

\begin{table}[t]
\centering
\caption{Maximum flux error against 500-point quadrature\label{tab:accuracy}}
\begin{tabular}{lc}
\hline\hline
Quadrature & Max.\ $|\Delta F|$ \\
\hline
$n_\mathrm{GL}=20$ & $3\times10^{-9}$ -- $1.4\times10^{-7}$ \\
$n_\mathrm{GL}=50$ (default) & $3\times10^{-11}$ -- $2.7\times10^{-10}$ \\
$n_\mathrm{GL}=100$ & $1\times10^{-12}$ -- $1.2\times10^{-11}$ \\
$n_\mathrm{GL}=200$ & $3\times10^{-14}$ -- $5.6\times10^{-13}$ \\
C++ default (20/50) & $3\times10^{-11}$ -- $3.6\times10^{-9}$ \\
\hline
\end{tabular}
\tablecomments{Ranges span the limb-darkening laws and $N_c = 0$--3. The
largest errors occur for planets just inside the stellar limb, where the
integrand has a square-root endpoint.}
\end{table}

\section{Performance} \label{sec:performance}

The GPUs available to us were shared with other jobs, so all GPU timings
here are conservative and indicative. Table~\ref{tab:timing} compares the
fixed-work model on an NVIDIA Quadro RTX~6000, which executes double
precision at 1/32 of its single-precision rate, with the C++ on one CPU
core. For $N_c=1$, the general model matches the speed of the earlier
first-order-only implementation from which it was derived. A value and
gradient for 500 points at $N_c=3$ takes 2.9\,ms on the GPU, against
14.4\,ms for the C++ primitive on the same node's CPU.

\begin{table}[t]
\centering
\caption{Time per light curve (ms)\label{tab:timing}}
\begin{tabular}{rcccc}
\hline\hline
Points & $N_c$ & C++, 1 core & GPU & GPU, value+grad \\
\hline
$10^3$ & 1 & 2.1 & 1.7 & 2.0 \\
$10^3$ & 2 & 4.8 & 2.3 & 3.2 \\
$10^3$ & 3 & 8.1 & 3.4 & 4.4 \\
$10^4$ & 1 & 22.0 & 16.0 & 17.8 \\
$10^4$ & 2 & 45.1 & 22.7 & 27.6 \\
$10^4$ & 3 & 77.7 & 29.6 & 28.0 \\
\hline
\end{tabular}
\tablecomments{GPU: a shared NVIDIA Quadro RTX~6000. The C++ column is
the flux only, without derivatives.}
\end{table}

For a workload like that of \citet{alderson2026}, order-six cosine-only
strings and 582 points at a 60\,s cadence, the C++ takes 19--20\,ms per
light curve on one core. The fixed-work model takes 3.2--3.8\,ms on the
RTX~6000, and on an NVIDIA V100, which has full-rate double precision,
1.7\,ms for one light curve and 0.81\,ms each in batches of 3120: about
24 CPU cores' worth. Large batches are evaluated in chunks of bounded
memory. Their cross-validated ensemble fit of one wavelength bin
($26\times5000\times10\times12 \approx 1.6\times10^7$ light curves) thus
needs ${\sim}85$ core-hours of C++ evaluations, or ${\sim}3.5$ hours on
one shared V100. As their reported cost per bin is about two weeks on 40
cores, most of it evidently lies outside the light-curve kernel.

\section{Issues found in the original implementation}
\label{sec:issues}

Validation turned up the following in the C++ and its \jax\ interface,
which v2's fixed-work model avoids. We have not changed the C++.
\begin{enumerate}
\item For planets that cover the stellar centre, a stellar arc crossing
the far side of the star loses $2\pi$ in its principal angle difference,
so the occulted integral can be wrong by a whole star. Fluxes above unity
result, e.g.\ $1.0003$ for an $a_0=1.5$ string whose true flux is
$0.0003$.
\item The half-angle conversion to the true anomaly is clamped to
$\sin f=0$, $\cos f=-1$ when $1+\cos E\leq10^{-10}$, which flattens the
orbit near apastron ($\partial\sin f/\partial M = 0$ instead of $-0.82$
at $e=0.1$).
\item The eccentric anomaly at transit uses the half-angle form, whose
derivatives with respect to $\omega$ are wrong near $\omega=-\pi/2$.
\item At $e=0$, $\partial F/\partial e$ does not depend on $\omega$
(e.g.\ $6\times10^{-4}$ against a true $0.086$).
\item The derivative of a circular planet's flux with respect to $a_0$
is in error by ${\sim}2\times10^{-6}$ just inside the second contact.
\item When an endpoint combination is absent from the segment lookup
table, no segment type is recorded, which misaligns the types of the
following segments (found by reading the code).
\item The \jax\ interface does not import on \jax\ $\geq0.6$, and it
broadcast each orbital and limb-darkening parameter to the times while
the C++ reads only the first element, silently truncating per-time
arrays. Both are fixed in v2.
\end{enumerate}

\section{Summary} \label{sec:summary}

The changes in \harmonica\ v2 are algorithmic, not a new wrapper around
the C++. Bracketing replaces eigenvalues for the intersections, the
implicit-function theorem supplies their derivatives, the occulted region
is assembled from the directions of the crossings, and the radial and
orbital formulae are rewritten to avoid cancellation. Together these
give every light-curve point the same static work, which XLA can compile
for GPUs, batch and differentiate. The release supports any $N_c$, the
quadratic, non-linear and power-2 laws, and selects the C++ automatically
on CPUs. The code is built on \jax\ \citep{jax2018}.

\begin{acknowledgments}
We thank D.\ Grant for developing \harmonica\ and making it public, and
an anonymous colleague for a careful review that found several of the
issues corrected here.
\end{acknowledgments}

\software{\harmonica\ \citep{grant2023}, \jax\ \citep{jax2018},
NumPyro \citep{phan2019}}

\bibliography{harmonica_v2}{}
\bibliographystyle{aasjournalv7}

\end{document}
